\documentclass[10pt,a4paper]{amsart}
\usepackage[margin=19mm]{geometry}
\usepackage{amsmath,amssymb,mathtools}
\usepackage[scr=rsfs,cal=euler]{mathalfa}
\usepackage{xfrac}
\usepackage[unicode,hidelinks,bookmarks=false]{hyperref}
\newcommand{\ie}{i.e.\ }
\newcommand{\cf}{cf.\ }
\newcommand{\resp}{respectively\ }
\newcommand{\Subsection}{\subsection}
\providecommand{\nobreakdash}{\nobreak\hbox{-}}
\setlength{\emergencystretch}{2em}
\multlinegap0pt
\usepackage{bm}
\usepackage{bbm}
\usepackage{dsfont}
\usepackage{xspace}
\usepackage{cancel}
\usepackage{mathtools}
\usepackage{amsthm}
\makeatletter
\@ifundefined{theorem}{\newtheorem{theorem}{Theorem}}{}
\@ifundefined{lemma}{\newtheorem{lemma}[theorem]{Lemma}}{}
\makeatother
\DeclareMathOperator{\sgn}{sgn}
\newcommand{\dimL}{\mathrm{dim}\,L_\lambda}
\newcommand{\su}{\mathfrak{su}}
\title[The high-dimension limit of characters of compact reductive ]{The high-dimension limit of characters of compact reductive Lie groups and restrictions on the production of quantum randomness}
\author{Piotr Borodako, Adam Sawicki}
\date{Source: arXiv:2510.24553v2 (CC BY 4.0)}
\begin{document}
\maketitle

\noindent\emph{Source and licence.} The high-dimension limit of characters of compact reductive Lie groups and restrictions on the production of quantum randomness. Version arXiv:2510.24553v2 (CC BY 4.0), distributed under the Creative Commons Attribution 4.0 International licence, as recorded in the arXiv licence metadata for this exact version and shown on the abstract page https://arxiv.org/abs/2510.24553v2. Full rights notice: licenses/arxiv-2510.24553-CC-BY-4.0.txt.\par\smallskip

\noindent\textit{Editorial scope.} Counted unit: the Theorem whose source label is \texttt{None} in the article (source0.tex lines 306--311, source label \texttt{None}), delivered together with the article's own complete original proof of it (source0.tex lines 313--389, one \texttt{proof} environment of 77 lines). No mathematics is written, repaired or re-derived: statement and proof are the article's own text line for line. Required context retained uncounted: eq:wcf (lines 249--252); eq:dimension\_formula (lines 254--257); lemma:weights (lines 409--416). Every definition, display and label of those lines is retained unchanged. Omitted from this package and not redistributed: 1--248, 253--253, 258--305, 312--312, 390--408, 417--806. Nothing inside a retained excerpt is omitted or altered: every retained body is delivered exactly as the article's lines are written, so no omission receipt of any kind accompanies this package. Citations: the retained excerpts carry no citation command at all, so no bibliography entry of the article is transcribed and no reference is redistributed. Reference adaptation: none was needed and none was made; every reference inside the delivered unit resolves inside it, so no unresolved reference or citation is produced. Printed slips kept literal: no printed word, symbol, hypothesis or number of the counted statement or of its proof was altered. Layout and class: the article's own class is \texttt{article} with packages including inputenc, amsmath, amssymb, amsthm, babel, geometry, lmodern, graphicx, authblk, hyperref, xcolor, nicematrix; the wrapper is the lane's pinned \texttt{amsart} editorial wrapper at 10pt on A4, which restates the theorem-like environments the retained text needs and adds the article's own macro definitions that the retained text needs (\texttt{\textbackslash dimL}, \texttt{\textbackslash sgn}, \texttt{\textbackslash su}), with extra wrapper packages (bm, bbm, dsfont, xspace, cancel, mathtools). The delivered numbering is the unit's own. Assets: no retained span contains a figure environment, a graphics-inclusion command, a PDF-page inclusion, a table or a TikZ picture; no image file is copied into this package, the package declares no asset dependency and no image file is redistributed with it. Licence: version arXiv:2510.24553v2 of this article is distributed under the Creative Commons Attribution 4.0 International licence, as recorded in the arXiv licence metadata for this exact version and shown on the abstract page https://arxiv.org/abs/2510.24553v2. A full rights notice is delivered at licenses/arxiv-2510.24553-CC-BY-4.0.txt. Characters: every retained excerpt is pure ASCII, so the delivered engine, XeLaTeX with its default Latin Modern text font, reports no missing character.
\begin{equation}
    \label{eq:wcf}
    \chi_\lambda(e^{ih}) = \frac{\sum_{w \in W} \sgn(w) e^{i(\lambda + \rho|w h)}}{\prod_{\alpha \in R_+} (e^{i(\alpha|h)/2} - e^{-i(\alpha|h)/2})}.
\end{equation}

\begin{equation}
    \label{eq:dimension_formula}
    \dim V_{\lambda} = \prod_{\alpha \in R_+} \frac{(\lambda + \rho|\alpha)}{(\rho|\alpha)}.
\end{equation}

\begin{lemma}
\label{lemma:weights}
An effective weight $\lambda'$, which fulfills
$$
\lambda' + \rho_b' = \pi_{bR_{deg}}(\lambda + \rho),
$$
where $\lambda$ is an integral weight, $\pi_{bR_{deg}}$ is the orthogonal projection onto the subspace spanned by the effective roots $bR_{deg}$, and $\rho_b'$ is the Weyl vector corresponding to the effective representation, belongs to the weight lattice of the effective representation.
\end{lemma}

\begin{theorem}
For any fixed singular element $g \neq I$ in SU$(3)$, corresponding to a Cartan element $h_0 \neq 0$, the normalized character vanishes in the high-dimension limit:
\[
\lim_{\lambda \to \infty} \frac{\chi_\lambda(g)}{\dimL} = 0.
\]
\end{theorem}

\begin{proof}
% We write the WCF as:
% \begin{equation}
%     \chi_\lambda(e^{ih}) = \frac{\sum_{w \in W} \sgn(w) e^{i(\eta|wh)}}{\prod_{\alpha \in R_+} \left(e^{i(\alpha|h)/2} - e^{-i(\alpha|h)/2}\right)},
% \end{equation}
% where $\eta := \lambda + \rho_{\su(3)}$ and we have used the Weyl-invariance of the Killing form to write the exponent as $e^{i(\eta|wh)}$. 


We begin with the Weyl Character Formula, Eq. \eqref{eq:wcf}. This formula holds for any regular element $h$. We will denote the positive roots of SU$(3)$ by $\alpha_1 = L_1 - L_2$, $\alpha_2 = L_2 - L_3$, and $\alpha_1 + \alpha_2 = L_1 - L_3$. The Weyl group is the permutation group of three elements, $S_3$. The Weyl vector is $\rho_{\su(3)} = \alpha_1 + \alpha_2 = (1,0,-1)$. $\alpha_1$ and $\alpha_2$ are the simple roots.

We now address the degenerate case, where the denominator is zero. This occurs when $h$ lies in a hyperplane orthogonal to some root $\alpha$. We consider a singular element $h_0$ such that $(\alpha_1|h_0) = 0$ but $(\alpha|h_0) \neq 0$ for $\alpha \in \{\alpha_2, \alpha_1+\alpha_2\}$. As our $\alpha_1 = L_1 - L_2$, this means exactly that $h_0$ has the following form:
\[
h_0 = \begin{pmatrix}
    a & 0 & 0 \\
    0 & a & 0 \\
    0 & 0 & b
\end{pmatrix},
\]
where $2a+b=0$ and $a \neq b$.

To evaluate the character at $h_0$, we consider a nearby regular point $h = h_0 + \delta$ and take the limit $\delta \to 0$. The WCF for $h$ reads:
\begin{equation}
    \chi_\lambda(e^{i(h_0+\delta)}) = \frac{1}{\prod_{\alpha \in \{\alpha_2, \alpha_1+\alpha_2\}} \left(e^{i(\alpha|h_0+\delta)/2} - e^{-i(\alpha|h_0+\delta)/2}\right)} \cdot \frac{\sum_{w \in W} \sgn(w) e^{i(\eta|w(h_0+\delta))}}{2i\sin\left(\frac{1}{2}(\alpha_1|\delta)\right)}
\end{equation}

where we have introduced $\eta := \lambda + \rho_{\su(3)}$. We can understand this degeneracy by looking at the action of the Weyl group. The degeneracy with respect to $\alpha_1$ means that the transposition of the first and the second element keeps $h_0$ the same. We can look at it from a geometrical point of view. The reflection of $h_0$ in the plane perpendicular to $\alpha_1$ is given by
\[
s_{\alpha_1}(h_0) = h_0 - \frac{2(h_0|\alpha_1)}{(\alpha_1|\alpha_1)}\alpha_1.
\]
As $(h_0|\alpha_1)$ is 0, we see that $h_0$ is a fixed point of this reflection. The subgroup of the Weyl group that leaves $h_0$ invariant is the stabilizer, $W_0 = \{e, s_{\alpha_1}\}$. We can isolate this stability by factoring the numerator sum over the cosets $b \in W/W_0$. A set of coset representatives is $\{e, s_{\alpha_1}, s_{\alpha_2}s_{\alpha_1}\}$. The factorization is:
\begin{equation}
    \sum_{w \in W} \sgn(w) e^{i(\eta|w(h_0+\delta))} = \sum_{b \in W/W_0} \sgn(b) e^{i(\eta|bh_0)} \left( \sum_{\sigma \in W_0} \sgn(\sigma) e^{i(\eta|b\sigma\delta)} \right).
\end{equation}
We analyze the singular part of the expression for each term $b$ by conjugating the stabilizer:
\begin{equation}
    \sum_{\sigma \in W_0} \sgn(\sigma) e^{i(\eta|b\sigma\delta)} = \sum_{\sigma \in bW_0b^{-1}} \sgn(\sigma) e^{i(\eta|\sigma b\delta)} = \sum_{\sigma \in \{e, s_{b\alpha_1}\}} \sgn(\sigma) e^{i(\eta|\sigma b\delta)}.
\end{equation}
If we assume that $\delta\in \text{span}(\alpha_1)$, we can interpret the ratio of this sum to the singular part of the denominator as a WCF for an SU$(2)$ sub-algebra given as a subspace $\mathfrak{g}_{b\alpha_1} \oplus \mathfrak{g}_{-b\alpha_1} \oplus [\mathfrak{g}_{b\alpha_1},\mathfrak{g}_{-b\alpha_1}] $  associated with the root $b\alpha_1$:
\begin{equation}
    \frac{\sum_{\sigma \in \{e, s_{b\alpha_1}\}} \sgn(\sigma) e^{i(\lambda + \rho_{\mathfrak{su}(3)}|\sigma b\delta)}}{2i\sin\left(\frac{1}{2}(b\alpha_1|\delta)\right)} = \frac{\sum_{\sigma \in \{e, s_{b\alpha_1}\}} \sgn(\sigma) e^{(\lambda'_b+\rho_{\su(2)}|\sigma b\delta)}}{2i\sin\left(\frac{1}{2}(b\alpha_1|\delta)\right)}
\end{equation}
where both exponents have to be equal, $\lambda'_b \in \text{span} (b\alpha_1)$ and $\rho_{\su(2)}=\frac{1}{2}b\alpha_1$. Thus the effective highest weight $\lambda'_{b}$\footnote{The fact that $\lambda'_b$ is a valid weight follows from the construction of the sub-representation, and will be proven explicitly in general in the next section in the lemma \ref{lemma:weights}.} of this sub-representation is defined by the condition:

\begin{equation}
    \label{eq:effective_weight}
    \lambda'_{b} + \rho_{\su(2)} = \frac{(\lambda + \rho_{\su(3)}|b\alpha_1)}{(b\alpha_1|b\alpha_1)} b\alpha_1.
\end{equation}
 In the limit $\delta \to 0$, this ratio gives the dimension of this sub-representation, $\dim L_{\lambda'_{b}}$:
\begin{equation}
    \frac{\sum_{\sigma \in \{e, s_{b\alpha_1}\}} \sgn(\sigma) e^{i(\eta|\sigma b\delta)}}{2i\sin\left(\frac{1}{2}(b\alpha_1|\delta)\right)} \xrightarrow{\delta \to 0} \dim L_{\lambda'_{b}}.
\end{equation}

Taking the limit of the full expression, we obtain the exact character at $h_0$:
\begin{equation}
    \chi_\lambda(h_0) = \frac{1}{\prod_{\alpha \in \{\alpha_2, \alpha_1+\alpha_2\}} \left(e^{i(\alpha|h_0)/2} - e^{-i(\alpha|h_0)/2}\right)} \sum_{b \in W/W_0} \sgn(b) e^{i(\eta|b h_0)} \dim L_{\lambda'_{b}}.
\end{equation}
To analyze the limit $||\lambda|| \to \infty$, we consider the ratio $|\chi_\lambda(h_0)|/\dimL$. By the triangle inequality, this is bounded by a finite sum of positive terms. The limit of this sum is zero if the limit of each term's ratio is zero. We therefore analyze the ratio for a representative term, $b=e$:

From Eq. (\ref{eq:dimension_formula}) we know that the exact dimension of the sub-representation is 
\begin{equation}
    \mathrm{dim}\,L_{\lambda'_{e}} = \frac{(\lambda' + \rho_{\su(2)}|\alpha_1)}{(\rho_{\su(2)}|\alpha_1)}= \frac{(\lambda + \rho_{\su(3)}|\alpha_1)}{(\rho_{\su(2)}|\alpha_1)}
\end{equation}

where we have used \ref{eq:effective_weight}. The ratio of this dimension to that of the full representation is:

\begin{align}
    \frac{\mathrm{dim}\,L_{\lambda'_{e}}}{\mathrm{dim}\,L_\lambda} 
    &= \frac{(\lambda + \rho_{\su(3)}|\alpha_1)}{(\rho_{\su(2)}|\alpha_1)} \left( \prod_{\alpha \in R_+} \frac{(\lambda+\rho_{\su(3)}|\alpha)}{(\rho_{\su(3)}|\alpha)} \right)^{-1} \nonumber \\[\jot]
    &= \frac{(\rho_{\su(3)}|\alpha_1)}{(\rho_{\su(2)}|\alpha_1)} \cdot \frac{(\rho_{\su(3)}|\alpha_2)(\rho_{\su(3)}|\alpha_1+\alpha_2)}{(\lambda+\rho_{\su(3)}|\alpha_2)(\lambda+\rho_{\su(3)}|\alpha_1+\alpha_2)}. \label{eq:ratio_limit_rigorous}
\end{align}

The numerator of this final expression in Eq. \eqref{eq:ratio_limit_rigorous} is a constant that depends only on the fixed Weyl 
vectors $\rho_{\su(3)}$ and $\rho_{\su(2)}$. The denominator is a product of terms that grow linearly with $\lambda$. 
Therefore, the ratio vanishes as $\lambda \to \infty$. Since every term in the bounding sum 
vanishes, the proof is complete.

\end{proof}

\end{document}
